\section{Separation axioms}
\label{section-separation-axioms}

\noindent
A topological space $X$ is Hausdorff if and only if the
diagonal $\Delta \subset X \times X$ is a closed subset.
The analogue in algebraic geometry is, given a scheme $X$ over
a base scheme $S$, to consider the
diagonal morphism
$$
\Delta_{X/S} : X \longrightarrow X \times_S X.
$$
This is the unique morphism of schemes such that
$\text{pr}_1 \circ \Delta_{X/S} = \text{id}_X$ and
$\text{pr}_2 \circ \Delta_{X/S} = \text{id}_X$ (it exists in
any category with fibre products).

\begin{lemma}
\label{lemma-diagonal-affines-closed}
The diagonal morphism of a morphism between affines is closed.
\end{lemma}

\begin{proof}
The diagonal morphism associated to the morphism
$\Spec(S) \to \Spec(R)$ is the morphism on spectra
corresponding to the ring
map $S \otimes_R S \to S$, $a \otimes b \mapsto ab$.
This map is clearly surjective, so $S \cong S \otimes_R S/J$
for some ideal $J \subset S \otimes_R S$. Hence
$\Delta$ is a closed immersion according to
Example \ref{example-closed-immersion-affines}.
\end{proof}

\begin{lemma}
\label{lemma-diagonal-immersion}
\begin{slogan}
The diagonal morphism for relative schemes is an immersion.
\end{slogan}
Let $X$ be a scheme over $S$.
The diagonal morphism $\Delta_{X/S}$ is an immersion.
\end{lemma}

\begin{proof}
Recall that if $V \subset X$ is affine open and maps into
$U \subset S$ affine open, then $V \times_U V$ is affine open
in $X \times_S X$, see Lemmas \ref{lemma-fibre-product-affines}
and \ref{lemma-open-fibre-product}.
Consider the open subscheme $W$ of $X \times_S X$ which
is the union of these affine opens $V \times_U V$.
By Lemma \ref{lemma-closed-local-target} it is enough
to show that each morphism
$\Delta_{X/S}^{-1}(V \times_U V) \to V \times_U V$ is
a closed immersion. Since $V = \Delta_{X/S}^{-1}(V \times_U V)$
we are just checking that $\Delta_{V/U}$ is a closed
immersion, which is Lemma \ref{lemma-diagonal-affines-closed}.
\end{proof}

\begin{definition}
\label{definition-separated}
Let $f : X \to S$ be a morphism of schemes.
\begin{enumerate}
\item We say $f$ is {\it separated} if the diagonal morphism $\Delta_{X/S}$
is a closed immersion.
\item We say $f$ is {\it quasi-separated} if the diagonal morphism
$\Delta_{X/S}$ is a quasi-compact morphism.
\item We say a scheme $Y$ is {\it separated} if the morphism
$Y \to \Spec(\mathbf{Z})$ is separated.
\item We say a scheme $Y$ is {\it quasi-separated} if the morphism
$Y \to \Spec(\mathbf{Z})$ is quasi-separated.
\end{enumerate}
\end{definition}

\noindent
By Lemmas \ref{lemma-diagonal-immersion} and \ref{lemma-immersion-when-closed}
we see that $\Delta_{X/S}$ is a closed immersion if an only if
$\Delta_{X/S}(X) \subset X \times_S X$ is a closed subset. Moreover,
by Lemma \ref{lemma-closed-immersion-quasi-compact} we see that a
separated morphism is quasi-separated. The reason for introducing
quasi-separated morphisms is that nonseparated morphisms come up naturally
in studying algebraic varieties (especially when doing moduli,
algebraic stacks, etc). But most often they are still quasi-separated.

\begin{example}
\label{example-not-quasi-separated}
Here is an example of a non-quasi-separated morphism.
Suppose $X = X_1 \cup X_2 \to S = \Spec(k)$ with
$X_1 = X_2 = \Spec(k[t_1, t_2, t_3, \ldots])$
glued along the complement of $\{0\} = \{(t_1, t_2, t_3, \ldots)\}$
(glued as in Example \ref{example-affine-space-zero-doubled}).
In this case the inverse image of the affine scheme
$X_1 \times_S X_2$ under $\Delta_{X/S}$ is the scheme
$\Spec(k[t_1, t_2, t_3, \ldots]) \setminus \{0\}$
which is not quasi-compact.
\end{example}

\begin{lemma}
\label{lemma-where-are-they-equal}
Let $X$, $Y$ be schemes over $S$.
Let $a, b : X \to Y$ be morphisms of schemes over $S$.
There exists a largest locally closed subscheme
$Z \subset X$ such that $a|_Z = b|_Z$. In fact $Z$ is
the equalizer of $(a, b)$. Moreover, if $Y$ is separated
over $S$, then $Z$ is a closed subscheme.
\end{lemma}

\begin{proof}
The equalizer of $(a, b)$ is for categorical reasons
the fibre product $Z$ in the following diagram
$$
\xymatrix{
Z = Y \times_{(Y \times_S Y)} X \ar[r] \ar[d] &
 X \ar[d]^{(a , b)} \\
Y \ar[r]^-{\Delta_{Y/S}} & Y \times_S Y
}
$$
Thus the lemma follows from Lemmas
\ref{lemma-base-change-immersion}, \ref{lemma-diagonal-immersion} and
Definition \ref{definition-separated}.
\end{proof}

\begin{lemma}
\label{lemma-characterize-quasi-separated}
Let $f : X \to S$ be a morphism of schemes.
The following are equivalent:
\begin{enumerate}
\item The morphism $f$ is quasi-separated.
\item For every pair of affine opens $U, V \subset X$
which map into a common affine open of $S$ the intersection
$U \cap V$ is a finite union of affine opens of $X$.
\item There exists an affine open covering $S = \bigcup_{i \in I} U_i$
and for each $i$ an affine open covering $f^{-1}U_i = \bigcup_{j \in I_i} V_j$
such that for each $i$ and each pair $j, j' \in I_i$ the
intersection $V_j \cap V_{j'}$ is a finite union of affine
opens of $X$.
\end{enumerate}
\end{lemma}

\begin{proof}
Let us prove that (3) implies (1).
By Lemma \ref{lemma-affine-covering-fibre-product}
the covering $X \times_S X = \bigcup_i \bigcup_{j, j'} V_j \times_{U_i} V_{j'}$
is an affine open covering of $X \times_S X$.
Moreover, $\Delta_{X/S}^{-1}(V_j \times_{U_i} V_{j'}) = V_j \cap V_{j'}$.
Hence the implication follows from Lemma \ref{lemma-quasi-compact-affine}.

\medskip\noindent
The implication (1) $\Rightarrow$ (2) follows from the fact
that under the hypotheses of (2) the fibre product
$U \times_S V$ is an affine open of $X \times_S X$.
The implication (2) $\Rightarrow$ (3) is trivial.
\end{proof}

\begin{lemma}
\label{lemma-characterize-separated}
Let $f : X \to S$ be a morphism of schemes.
\begin{enumerate}
\item If $f$ is separated then for every pair of affine
opens $(U, V)$ of $X$ which map into a
common affine open of $S$ we have
\begin{enumerate}
\item the intersection $U \cap V$ is affine.
\item the ring map
$\mathcal{O}_X(U) \otimes_{\mathbf{Z}} \mathcal{O}_X(V)
\to \mathcal{O}_X(U \cap V)$
is surjective.
\end{enumerate}
\item If any pair of points $x_1, x_2 \in X$ lying over a common
point $s \in S$ are contained in affine opens $x_1 \in U$,
$x_2 \in V$ which map into a common affine open of $S$ such
that (a), (b) hold, then $f$ is separated.
\end{enumerate}
\end{lemma}

\begin{proof}
Assume $f$ separated. Suppose $(U, V)$ is a pair as in (1).
Let $W = \Spec(R)$ be an affine open of $S$ containing
both $f(U)$ and $f(V)$. Write $U = \Spec(A)$ and
$V = \Spec(B)$ for $R$-algebras $A$ and $B$.
By Lemma \ref{lemma-open-fibre-product} we see that
$U \times_S V = U \times_W V = \Spec(A \otimes_R B)$
is an affine open of $X \times_S X$. Hence, by
Lemma \ref{lemma-closed-subspace-scheme} we see that
$\Delta^{-1}(U \times_S V) \to U \times_S V$
can be identified with $\Spec((A \otimes_R B)/J) \to \Spec(A \otimes_R B)$
for some ideal $J \subset A \otimes_R B$.
Thus $U \cap V = \Delta^{-1}(U \times_S V)$ is affine.
Assertion (1)(b) holds because
$A \otimes_{\mathbf{Z}} B \to (A \otimes_R B)/J$ is surjective.

\medskip\noindent
Assume the hypothesis formulated in (2) holds.
Clearly the collection of affine opens $U \times_S V$
for pairs $(U, V)$ as in (2) form an affine open covering
of $X \times_S X$ (see e.g.\ Lemma \ref{lemma-affine-covering-fibre-product}).
Hence it suffices to show that each morphism
$U \cap V = \Delta_{X/S}^{-1}(U \times_S V) \to U \times_S V$
is a closed immersion, see Lemma \ref{lemma-closed-local-target}.
By assumption (a) we have $U \cap V = \Spec(C)$ for some ring $C$.
After choosing an affine open $W = \Spec(R)$ of $S$
into which both $U$ and $V$ map and writing $U = \Spec(A)$,
$V = \Spec(B)$ we see that the assumption (b) means
that the composition
$$
A \otimes_{\mathbf{Z}} B \to A \otimes_R B \to C
$$
is surjective. Hence $A \otimes_R B \to C$ is surjective and
we conclude that $\Spec(C) \to \Spec(A \otimes_R B)$
is a closed immersion.
\end{proof}

\begin{example}
\label{example-projective-line-separated}
Let $k$ be a field. Consider the structure morphism
$p : \mathbf{P}^1_k \to \Spec(k)$ of the projective
line over $k$, see Example \ref{example-projective-line}.
Let us use the lemma above to prove that $p$
is separated. By construction $\mathbf{P}^1_k$ is covered by two
affine opens $U = \Spec(k[x])$ and $V = \Spec(k[y])$
with intersection $U \cap V = \Spec(k[x, y]/(xy - 1))$
(using obvious notation). Thus it suffices to check that
conditions (2)(a) and (2)(b) of Lemma \ref{lemma-characterize-separated}
hold for the pairs of affine opens $(U, U)$, $(U, V)$, $(V, U)$
and $(V, V)$. For the pairs $(U, U)$ and $(V, V)$ this is trivial.
For the pair $(U, V)$ this amounts to proving
that $U \cap V$ is affine, which is true, and that the ring map
$$
k[x] \otimes_{\mathbf{Z}} k[y] \longrightarrow k[x, y]/(xy - 1)
$$
is surjective. This is clear because any element in the
right hand side can be written as a sum of a polynomial
in $x$ and a polynomial in $y$.
\end{example}

\begin{lemma}
\label{lemma-fibre-product-after-map}
Let $f : X \to T$ and $g : Y \to T$ be morphisms of schemes
with the same target. Let $h : T \to S$ be a morphism of schemes.
Then the induced morphism $i : X \times_T Y \to X \times_S Y$ is
an immersion. If $T \to S$ is separated, then $i$ is a closed
immersion. If $T \to S$ is quasi-separated, then $i$ is a
quasi-compact morphism.
\end{lemma}

\begin{proof}
By general category theory the following diagram
$$
\xymatrix{
X \times_T Y \ar[r] \ar[d] & X \times_S Y \ar[d] \\
T \ar[r]^{\Delta_{T/S}} \ar[r] & T \times_S T
}
$$
is a fibre product diagram. The lemma follows
from Lemmas \ref{lemma-diagonal-immersion},
\ref{lemma-fibre-product-immersion} and
\ref{lemma-quasi-compact-preserved-base-change}.
\end{proof}

\begin{lemma}
\label{lemma-semi-diagonal}
Let $g : X \to Y$ be a morphism of schemes over $S$.
The morphism $i : X \to X \times_S Y$ is an immersion.
If $Y$ is separated over $S$ it is a closed immersion.
If $Y$ is quasi-separated over $S$ it is quasi-compact.
\end{lemma}

\begin{proof}
This is a special case of Lemma \ref{lemma-fibre-product-after-map}
applied to the morphism $X = X \times_Y Y \to X \times_S Y$.
\end{proof}

\begin{lemma}
\label{lemma-section-immersion}
Let $f : X \to S$ be a morphism of schemes.
Let $s : S \to X$ be a section of $f$ (in a formula $f \circ s = \text{id}_S$).
Then $s$ is an immersion.
If $f$ is separated then $s$ is a closed immersion.
If $f$ is quasi-separated, then $s$ is quasi-compact.
\end{lemma}

\begin{proof}
This is a special case of Lemma \ref{lemma-semi-diagonal} applied to
$g =s$ so the morphism $i = s : S \to S \times_S X$.
\end{proof}

\begin{lemma}
\label{lemma-separated-permanence}
Permanence properties.
\begin{enumerate}
\item A composition of separated morphisms is separated.
\item A composition of quasi-separated morphisms is quasi-separated.
\item The base change of a separated morphism is separated.
\item The base change of a quasi-separated morphism is quasi-separated.
\item A (fibre) product of separated morphisms is separated.
\item A (fibre) product of quasi-separated morphisms is quasi-separated.
\end{enumerate}
\end{lemma}

\begin{proof}
Let $X \to Y \to Z$ be morphisms. Assume that $X \to Y$ and
$Y \to Z$ are separated. The composition
$$
X \to X \times_Y X \to X \times_Z X
$$
is closed because the first one is by assumption and the second
one by Lemma \ref{lemma-fibre-product-after-map}. The same argument
works for ``quasi-separated'' (with the same references).

\medskip\noindent
Let $f : X \to Y$ be a morphism of schemes over a base $S$.
Let $S' \to S$ be a morphism of schemes. Let $f' : X_{S'} \to Y_{S'}$
be the base change of $f$. Then the diagonal morphism
of $f'$ is a morphism
$$
\Delta_{f'} :
X_{S'} = S' \times_S X
\longrightarrow
X_{S'} \times_{Y_{S'}} X_{S'} = S' \times _S (X \times_Y X)
$$
which is easily seen to be the base change of $\Delta_f$.
Thus (3) and (4) follow from the fact that
closed immersions and quasi-compact morphisms are preserved
under arbitrary base change (Lemmas
\ref{lemma-fibre-product-immersion} and
\ref{lemma-quasi-compact-preserved-base-change}).

\medskip\noindent
If $f : X \to Y$ and $g : U \to V$ are morphisms of schemes over a base $S$,
then $f \times g$ is the composition of $X \times_S U \to X \times_S V$
(a base change of $g$) and $X \times_S V \to Y \times_S V$ (a base change
of $f$). Hence (5) and (6) follow from (1) -- (4).
\end{proof}

\begin{lemma}
\label{lemma-compose-after-separated}
\begin{slogan}
Separated and quasi-separated morphisms satisfy cancellation.
\end{slogan}
Let $f : X \to Y$ and $g : Y \to Z$ be morphisms of schemes.
If $g \circ f$ is separated then so is $f$.
If $g \circ f$ is quasi-separated then so is $f$.
\end{lemma}

\begin{proof}
Assume that $g \circ f$ is separated.
Consider the factorization $X \to X \times_Y X \to X \times_Z X$
of the diagonal morphism of $g \circ f$.
By Lemma \ref{lemma-fibre-product-after-map}
the last morphism is an immersion. By assumption the image
of $X$ in $X \times_Z X$ is closed. Hence it is also closed
in $X \times_Y X$. Thus we see that $X \to X \times_Y X$
is a closed immersion by Lemma \ref{lemma-immersion-when-closed}.

\medskip\noindent
Assume that $g \circ f$ is quasi-separated.
Let $V \subset Y$ be an affine open which maps into an affine
open of $Z$. Let $U_1, U_2 \subset X$ be affine opens which
map into $V$. Then $U_1 \cap U_2$ is a finite union of affine
opens because $U_1, U_2$ map into a common affine open
of $Z$. Since we may cover $Y$ by affine opens like $V$ we
deduce the lemma from Lemma \ref{lemma-characterize-quasi-separated}.
\end{proof}

\begin{lemma}
\label{lemma-quasi-compact-permanence}
Let $f : X \to Y$ and $g : Y \to Z$ be morphisms of schemes.
If $g \circ f$ is quasi-compact and $g$ is quasi-separated
then $f$ is quasi-compact.
\end{lemma}

\begin{proof}
This is true because $f$ equals the composition
$(1, f) : X \to X \times_Z Y \to Y$. The first map
is quasi-compact by Lemma \ref{lemma-section-immersion}
because it is a section of the quasi-separated morphism $X \times_Z Y \to X$
(a base change of $g$, see Lemma \ref{lemma-separated-permanence}).
The second map is quasi-compact as it
is the base change of $g \circ f$, see
Lemma \ref{lemma-quasi-compact-preserved-base-change}.
And compositions of quasi-compact
morphisms are quasi-compact, see Lemma \ref{lemma-composition-quasi-compact}.
\end{proof}

\begin{lemma}
\label{lemma-affine-separated}
An affine scheme is separated. A morphism from an affine scheme
to another scheme is separated.
\end{lemma}

\begin{proof}
Let $U = \Spec(A)$ be an affine scheme. Then $U \to \Spec(\mathbf{Z})$
has closed diagonal by Lemma \ref{lemma-diagonal-affines-closed}.
Thus $U$ is separated by Definition \ref{definition-separated}.
If $U \to X$ is a morphism of schemes, then we can apply
Lemma \ref{lemma-compose-after-separated}
to the morphisms $U \to X \to \Spec(\mathbf{Z})$ to conclude
that $U \to X$ is separated.
\end{proof}

\noindent
You may have been wondering whether the condition
of only considering pairs of affine opens whose image is contained
in an affine open is really necessary to be able to conclude
that their intersection is affine. Often it isn't!

\begin{lemma}
\label{lemma-curiosity}
Let $f : X \to S$ be a morphism.
Assume $f$ is separated and $S$ is a separated scheme.
Suppose $U \subset X$ and $V \subset X$ are affine open.
Then $U \cap V$ is affine (and a closed subscheme of $U \times V$).
\end{lemma}

\begin{proof}
In this case $X$ is separated by Lemma \ref{lemma-separated-permanence}.
Hence $U \cap V$ is affine by
applying Lemma \ref{lemma-characterize-separated} to the
morphism $X \to \Spec(\mathbf{Z})$.
\end{proof}

\noindent
On the other hand, the following example shows that we cannot
expect the image of an affine to be contained in an affine.

\begin{example}
\label{example-image-affine-projective}
Consider the nonaffine scheme
$U = \Spec(k[x, y]) \setminus \{(x, y)\}$ of
Example \ref{example-not-affine}. On the other hand, consider the
scheme
$$
\mathbf{GL}_{2, k} = \Spec(k[a, b, c, d, 1/ad - bc]).
$$
There is a morphism $\mathbf{GL}_{2, k} \to U$ corresponding
to the ring map $x \mapsto a$, $y \mapsto b$. It is easy to see that
this is a surjective morphism, and hence the image is not contained
in any affine open of $U$. In fact, the affine scheme
$\mathbf{GL}_{2, k}$ also surjects onto $\mathbf{P}^1_k$, and
$\mathbf{P}^1_k$ does not even have an immersion into {\it any} affine scheme.
\end{example}

\begin{remark}
\label{remark-quasi-compact-and-quasi-separated}
Let $P$ equal one of the following $4$ properties:
``quasi-separated'', ``separated'',
``quasi-compact and quasi-separated'', or ``quasi-compact and separated''.
Then the following are true
\begin{enumerate}
\item Any affine scheme has $P$.
\item Given a morphism $f : X \to Y$ of scheme where $Y$ has $P$, then
$f$ has $P$ if and only if $X$ has $P$.
\item If  $X \to Y$ and $Z \to Y$ are morphisms of schemes
and $X$, $Y$, $Z$ have $P$, then $X \times_Y Z$ has $P$.
\end{enumerate}
Statement (1) is clear.
Let $f : X \to Y$ be as in (2). If $f$ has $P$, then so does $X$
by Lemmas \ref{lemma-separated-permanence} and
\ref{lemma-composition-quasi-compact} applied to
$X \to Y \to \Spec(\mathbf{Z})$. If $X$ has $P$, then so does $f$ by
Lemmas \ref{lemma-compose-after-separated} and
\ref{lemma-quasi-compact-permanence} applied to
$X \to Y \to \Spec(\mathbf{Z})$.
Let $X \to Y$ and $Z \to Y$ be as in (3).
Then the projection $X \times_Y Z \to Z$ has $P$ as a base change of
the arrow $X \to Y$ which has $P$ by (2), see Lemmas
\ref{lemma-separated-permanence} and
\ref{lemma-quasi-compact-preserved-base-change}.
Hence $X \times_Y Z$ has $P$ by (2).
\end{remark}





